% Part: computability
% Chapter: recursive-functions
% Section: trees

\documentclass[../../../include/open-logic-section]{subfiles}

\begin{document}

\olfileid[hi]{cmp}{rec}{tre}
\olsection{वृक्ष}

कभी-कभी वृक्षों को प्राकृतिक संख्याओं से निरूपित करना उपयोगी होता है,
ठीक वैसे ही जैसे हम अनुक्रमों को संख्याओं से और उनके गुणों तथा उन पर होने
वाली संक्रियाओं को उनके कूटों पर आदिम पुनरावर्ती संबंधों और फलनों से
निरूपित कर सकते हैं। ऐसा करने के लिए हम अनुक्रमों और उनके कूटों का उपयोग
करेंगे। कोई वृक्ष या तो एक अकेला शीर्ष (संभवतः किसी चिह्नक सहित) हो सकता
है, या अनेक उपवृक्षों से जुड़ा कोई शीर्ष (संभवतः किसी चिह्नक सहित)। उस
शीर्ष को वृक्ष का \emph{मूल} और उससे जुड़े उपवृक्षों को उसके
\emph{तत्क्षण उपवृक्ष} कहते हैं।

हम वृक्षों को पुनरावर्ती रूप से अनुक्रम $\tuple{k, d_1, \dots, d_k}$
के रूप में कूटबद्ध करते हैं, जहाँ $k$ तत्क्षण उपवृक्षों की संख्या है और
$d_1$, \dots,~$d_k$ उन तत्क्षण उपवृक्षों के कूट हैं। यदि शीर्षों पर
चिह्नक हों, तो उन्हें तत्क्षण उपवृक्षों के बाद सम्मिलित किया जा सकता है।
अतः केवल चिह्नक~$l$ वाले एक शीर्ष से बने वृक्ष का कूट $\tuple{0,l}$
होगा; और ऐसे मूल (जिसका चिह्नक~$l_1$ है) से बने वृक्ष का कूट, जो दो
अकेले शीर्षों (जिनके चिह्नक $l_2$, $l_3$ हैं) से जुड़ा हो,
$\tuple{2,
  \tuple{0, l_2}, \tuple{0, l_3}, l_1}$ होगा।

\begin{prop}
  \ollabel{prop:subtreeseq}
  फलन $\fn{SubtreeSeq}(t)$, जो ऐसे अनुक्रम का कूट लौटाता है जिसके अवयव
  कूट~$t$ वाले वृक्ष के सभी उपवृक्षों के कूट हैं, आदिम पुनरावर्ती है।
\end{prop}

\begin{proof}
  पहले ध्यान दें कि $\fn{ISubtrees}(t) = \fn{subseq}(t, 1, (t)_0)$
  आदिम पुनरावर्ती है और वृक्ष~$t$ के तत्क्षण उपवृक्षों के कूट लौटाता है।
  अब हम एक सहायक फलन $\fn{hSubtreeSeq}(t,n)$ परिभाषित कर सकते हैं, जो
  मूल से $n$ किनारियों की दूरी पर स्थित सभी उपवृक्षों का अनुक्रम संगणित करता है।
  ~$t$ के उपवृक्षों का वह अनुक्रम जो मूल से $0$ किनारियों की दूरी पर है---दूसरे शब्दों
  में,~$t$ के मूल से आरंभ होने वाला अनुक्रम---केवल~$t$ से बना अनुक्रम
  है। स्तर~$n+1$ पर स्थित, $t$ के सभी उपवृक्षों का अनुक्रम पाने के लिए हम स्तर
  $n$ के उपवृक्षों का संलग्नन ऐसे अनुक्रम से करते हैं जो स्तर $n$ के
  उपवृक्षों के सभी तत्क्षण उपवृक्षों से बना है। इन सभी की सूची पाने के
  लिए ध्यान दें कि यदि $f(x)$ अनुक्रमों के कूट लौटाने वाला आदिम
  पुनरावर्ती फलन है, तो $g_f(s, k) = f((s)_0) \concat
  \dots \concat f((s)_k)$ भी आदिम पुनरावर्ती है:
    \begin{align*}
      g(s, 0) & = f((s)_0)\\
      g(s, k+1) & = g(s, k) \concat f((s)_{k+1})
    \end{align*}
    मसलन, यदि $s$ वृक्षों का अनुक्रम है, तो
    $h(s) = g_{\fn{ISubtrees}}(s, \len{s})$,~$s$ के अवयवों के तत्क्षण
    उपवृक्षों का अनुक्रम देता है। इसका उपयोग करके हम
    $\fn{hSubtreeSeq}$ को इस प्रकार परिभाषित कर सकते हैं:
    \begin{align*}
      \fn{hSubtreeSeq}(t, 0) & = \tuple{t} \\
      \fn{hSubtreeSeq}(t, n+1) & = \fn{hSubtreeSeq}(t, n) \concat
      h(\fn{hSubtreeSeq}(t, n)).
    \end{align*}
    कूट~$t$ वाले वृक्ष में उपवृक्षों का महत्तम स्तर, अर्थात् मूल और किसी
    पर्ण-शीर्ष के बीच की महत्तम दूरी, कूट~$t$ से परिबद्ध है। अतः कूट~$t$
    वाले वृक्ष के सभी उपवृक्षों के कूटों का अनुक्रम
    $\fn{hSubtreeSeq}(t, t)$ से मिलता है।
\end{proof}

\begin{prob}
  \olref[cmp][rec][tre]{prop:subtreeseq} के प्रमाण में
  $\fn{hSubtreeSeq}$ की परिभाषा सामान्यतः पुनरावृत्तियाँ सम्मिलित करती
  है। ऐसी वैकल्पिक परिभाषा दें जो सुनिश्चित करे कि परिणामी सूची में किसी
  उपवृक्ष का कूट केवल एक बार आए।
\end{prob}

\end{document}
